\section{Conclusions and reproducibility boundary}\label{sec:conclusion}
A bounded, clean-source retrospective study shows 4.4--5.9\% lower pooled MAE than contemporaneous persistence across four horizons and two cleaning regimes. MASE remains approximately 1.00--1.11 using training-period scales, with 4.4--4.5\% of rows omitted only from that metric. The conclusion is modest historical average error reduction under specified settings, not universal predictive superiority or improved farmer profit.

The contribution is a reproducible small-budget cleaning/dependence sensitivity study on retained inputs. Historical national-scale comparisons and the source-completeness audit are separated as supporting material; their original full-size execution and panel selection are not certified by this work. Dependence sensitivities remain approximate and retrospective. Frozen-model future issuance can begin collecting a different, explicitly defined population, but no future confirmation or human statistical approval is claimed.

\paragraph{Artifacts and availability.}
Fitting source is tagged \texttt{agdata-}\allowbreak\texttt{research-}\allowbreak\texttt{source-v1}; release records bind pre/post source/input/JAR hashes, commands and feature outputs. Saved predictions and generated primary tables are recorded under \path{results/release_v1}. Models, large data and issuance batches remain local outside Git. Public source/results tags and compact manifests do not constitute a public immutable data archive, independently witnessed execution or anonymised journal submission package. Later API responses can differ from retained bytes after revisions. The fixed panel is a declared external input. Corresponding-author details, declarations and formal journal-template compliance remain unverified.

The evidence ledger, commands, third-review response and future protocol distinguish reproduced retrospective results, historical supporting evidence and pending forward evaluation. Human expert statistical review, independent archival/reproduction and completed prospective outcomes remain necessary before a stronger submission recommendation.
